\section{Лекция 16}

\href{https://github.com/dotnet/BenchmarkDotNet}{BenchmarkDotNet} --- автоматизация замеров производительности для .NET.

\href{https://github.com/AndreyAkinshin/perfolizer}{Perfolizer} --- набор вспомогательных инструментов для анализа производительности.

Примеры на F\#.
\begin{itemize}
    \item \href{https://github.com/YaccConstructor/GraphBLAS-sharp/tree/dev/benchmarks/GraphBLAS-sharp.Benchmarks}{Корневая папка с бенчмарками}. 
    \item \href{https://github.com/YaccConstructor/GraphBLAS-sharp/blob/dev/benchmarks/GraphBLAS-sharp.Benchmarks/Program.fs}{Точка входа.} Указываем, какие бенчмарки доступны.
    \item \href{https://github.com/YaccConstructor/GraphBLAS-sharp/blob/dev/benchmarks/GraphBLAS-sharp.Benchmarks/BenchmarksEWiseAdd.fs}{Пример конкретных бенчмарков.}
\end{itemize}


Actor-ориентированное программирование как реализация асинхронного (concurrent) программирования. Коммуникация на сообщениях. Mailbox processor и Hopac как реализации. Примеры использования Mailbox processor и Hopac.

Домашняя работа 5

